\documentclass[11pt]{article}
\usepackage[margin=1in]{geometry}
\usepackage{amsmath,amssymb,booktabs,hyperref,graphicx,microtype}
\hypersetup{colorlinks=true,citecolor=blue,urlcolor=blue,linkcolor=blue}
\newcommand{\E}{\mathbb{E}}
\newcommand{\D}{\Delta}

\title{Curvature of Gradient-Boosted Trees:\\
Main Effects as Bending, Interactions as Off-Diagonal Curvature,\\
and Depth as Curvature Order}
\author{Nupur Krishnan \and Agus Sudjianto}
\date{\today}

\begin{document}
\maketitle

\begin{abstract}
Curvature is a second-order notion, and for a fitted predictor its second-order
structure splits in two: how each feature's effect bends along its own axis, and
how two features bend jointly. The second is the interaction. We define both for
gradient-boosted decision trees (GBDTs), which are piecewise constant and have no
classical Hessian, using finite-difference operators against a background
distribution. A main effect is a first difference, an interaction is a mixed
second difference and an order-$|S|$ effect is an $|S|$-fold difference;
interactions are the off-diagonal entries of this discrete Hessian. The same
object is read from the model-induced kernel: the component kernel's feature map
is the mixed-difference response, its diagonal $K_S(x,x)$ is a per-point
interaction-energy field and an additive model has a separable kernel. Tree depth
bounds the order of mixed difference a model can carry, so depth corresponds to
interaction order. On simulated data the kernel detects a planted interaction at
$267\times$ the energy of a permitted-but-absent one, recovers its shape at
Spearman $0.98$ and localizes where it fires at $0.98$, and each planted order is
captured when depth reaches it. On California housing and bikeshare an
importance matrix and a Hessian matrix, sharing the interaction off-diagonal,
identify the interacting pairs and separate important-but-linear main effects from
curved ones, and the depth-order curve places California's effective interaction
order near two and bikeshare's near three.
\end{abstract}

\medskip
\noindent\textbf{Keywords:} gradient boosting; functional ANOVA; interactions;
curvature; discrete Hessian; model-induced kernel; interaction detection; tree depth.

\section{Introduction}

A fitted model's first-order structure is its gradient: how the prediction moves
as one feature changes. Its second-order structure is the Hessian, and the
Hessian has two parts. The diagonal $\partial^2 f/\partial x_j^2$ is how the
effect of $x_j$ bends along its own axis. The off-diagonal
$\partial^2 f/\partial x_j\partial x_k$ is how the slope in $x_j$ changes as $x_k$
moves, which is the interaction between $x_j$ and $x_k$. A model is additive
exactly when every off-diagonal vanishes, so an interaction is nonzero
off-diagonal curvature.

Gradient-boosted trees complicate this because they are piecewise constant: the
classical Hessian is zero inside a leaf and undefined at split faces. We replace
derivatives with finite differences against a background distribution and recover
the same split. A main effect is a first difference in one coordinate; an
interaction is a mixed second difference in two; an order-$|S|$ effect is the
$|S|$-fold mixed difference. This is the functional ANOVA (fANOVA) reading of tree
ensembles \cite{Sudjianto2025,Hu2023}, read here as a curvature statement.

Three results follow. First, the interaction is visible in the model-induced
kernel: the component kernel's feature map is the mixed-difference response, its
diagonal is a per-point interaction-energy field and an additive model has a
kernel that separates into a sum over single features \cite{SudjiantoZhang2026}.
Second, tree depth bounds the order of mixed difference the model can represent,
so depth corresponds to interaction order rather than to local sharpness. Third,
both statements are testable, and we test them on simulated data with known
structure and on California housing and bikeshare.

\section{Curvature of a fitted model}

\subsection{The Hessian and its two parts}

For a twice-differentiable $f$, the Hessian $H_{jk}(x)=\partial^2 f/\partial
x_j\partial x_k$ carries the second-order behavior. The diagonal $H_{jj}$ is the
self-curvature of feature $j$: how sharply its effect bends. The off-diagonal
$H_{jk}$ ($j\neq k$) is the interaction: the part of the curvature that couples
two directions, the twist of the surface. Additivity is the statement that $H$ is
diagonal everywhere.

\subsection{Discrete curvature for piecewise-constant models}

A tree ensemble has no usable Hessian, so fix a background measure with draws
$B^{(1)},\dots,B^{(M)}$ and define difference operators. The main-effect (first
difference) and interaction (mixed second difference) contrasts are
\begin{align}
 \Psi_j(x)   &= \E_b\!\big[f(x_j,b_{-j})-f(b)\big],\\
 \Psi_{jk}(x)&= \E_b\!\big[f(x_j,x_k,b_{-jk})-f(x_j,b_{-j})-f(x_k,b_{-k})+f(b)\big].
\label{eq:contrasts}
\end{align}
$\Psi_j$ is the discrete $\partial_j f$ and $\Psi_{jk}$ is the discrete
$\partial^2_{jk}f$: the inclusion--exclusion in \eqref{eq:contrasts} is exactly
the mixed second difference $\D_j\D_k f$. The general order-$|S|$ contrast is the
$|S|$-fold difference
\begin{equation}
 \Psi_S(x)=\E_b\Big[\textstyle\sum_{T\subseteq S}(-1)^{|S|-|T|}\,f(x_T,b_{-T})\Big],
\label{eq:contrastS}
\end{equation}
the discrete $|S|$-th mixed derivative. The fANOVA decomposition of $f$ is
therefore a grading by the order of its mixed-difference tensor: gradient gives
main effects, the off-diagonal Hessian gives pairwise interactions, the
third-order tensor gives triple interactions. Fitting the deep component on the
main-effect residual makes $\Psi_S$ orthogonal to lower orders under the
background measure \cite{Hu2023}, so the estimate is the pure order-$|S|$ part.

\subsection{Curvature energy and two matrices}

The energy of a difference is $A_S=\E[(\E_b\Psi_S)^2]=\mathrm{Var}(g_S)$ with
$g_S=\E_b\Psi_S$ centered. Two second-order objects organize the model. The
interaction between $j$ and $k$ is the mixed second difference $\Psi_{jk}$, with
energy $A_{jk}$: the off-diagonal curvature coupling the two features. The
self-curvature of feature $j$ is the self second difference
\begin{equation}
 \Psi_{jj}(x)=\E_b\!\big[f(x_j{+}h,b_{-j})-2f(x_j,b_{-j})+f(x_j{-}h,b_{-j})\big],
\end{equation}
the discrete $\partial^2 f/\partial x_j^2$, with energy $A_{jj}$: how much the
main effect bends. A two-state feature has no interior step, so $A_{jj}=0$; its
main effect is a level shift, not a bend.

These give two matrices with a shared off-diagonal. The Hessian matrix
$H_{jk}=A_{jk}$, $H_{jj}=A_{jj}$ has a second difference in every entry, so it is
consistent in derivative order, and additive $f$ makes it diagonal. The importance
matrix $C_{jk}=A_{jk}$, $C_{jj}=A_j$ keeps the interaction off-diagonal but records
the first-difference energy of the main effect on the diagonal. Both are
consistent with the model-induced kernel: every entry is the mean local energy
$\E[K_S(x,x)]$ of the corresponding component's kernel (Section~\ref{sec:kernel}).
They differ only in whether the diagonal reads a main effect's importance or its
bending, so a feature can be important and linear (large $C_{jj}$, small $H_{jj}$)
or important and nonlinear.

\section{The model-induced kernel view}
\label{sec:kernel}

\subsection{Leaf co-membership kernel}

A fitted tree induces a kernel. Route $x$ to its leaf $\ell(x)$, write the one-hot
$z(x)=e_{\ell(x)}$, and $K_T(x,x')=z(x)^\top z(x')=\mathbf 1\{\ell(x)=\ell(x')\}$
is positive semidefinite since the Gram matrix is $ZZ^\top$. A GBDT sums these, so
two points are close when the ensemble routes them together
\cite{SudjiantoZhang2026}.

\subsection{Component kernel and non-separability}

The interaction contrast supplies a feature map $\Psi_S(x)\in\mathbb R^M$ and a PSD
kernel $K_S(x,x')=\tfrac1M\Psi_S(x)^\top\Psi_S(x')$: the same recipe with the
model's mixed-difference response in place of leaf indicators. Because each deep
tree is constrained to one pair, its leaf map depends jointly on $(x_j,x_k)$ and
does not separate, so the model kernel decomposes as $K=\sum_S K_S$ and the
off-diagonal blocks are the interactions. In kernel terms an interaction is a
non-separable block: an additive model has $K=\sum_j K_j$, and the mixed-derivative
operator $\D_j\D_k$ corresponds to the product kernel $k_j\,k_k$, the order-$|S|$
term of the ANOVA-kernel expansion $\prod_j(1+k_j)=\sum_S\prod_{j\in S}k_j$.

\subsection{Local interaction-energy field}

The diagonal of the component Gram matrix is a per-point field,
\begin{equation}
 K_S(x,x)=\tfrac1M\lVert\Psi_S(x)\rVert^2,
\end{equation}
the local squared mixed curvature: where in input space the interaction fires. Its
average is the importance $A_S$. So the kernel gives both a scalar strength and a
map of where the interaction is active.

\section{Depth and the order of curvature}

A depth-$d$ tree partitions on at most $d$ features along any root-to-leaf path, so
its output is a sum of terms each depending on at most $d$ features. Sums of such
trees stay order-$d$. Depth therefore bounds the order of mixed difference the
model can carry: a depth-$d$ ensemble represents interactions up to order $d$ and
none higher. Depth buys interaction order, not local sharpness, which is set
instead by the number of splits and boosting rounds at a fixed order. The clean
prediction is that an order-$m$ effect is captured when depth first reaches $m$.

\section{Experiments}

\subsection{Protocol}

Positive responses are transformed by $y=\log(1+\text{target})$. Contrasts use $96$
background draws. On the real data the effect-energy matrix and interaction fields
come from an unconstrained depth-six XGBoost; the depth-order curve fits
unconstrained models at each depth and scores held-out mean squared error. Models
run under XGBoost 3.0.

\subsection{Simulated: detection, recovery and localization}

We draw five independent uniform features, plant one localized interaction on
$(u_1,u_2)$ (a Gaussian bump in a corner) and leave $(u_3,u_4)$ additive, then fit
the staged model with both pairs permitted. Table~\ref{tab:sim-kernel} reports the
kernel view against the true purified interaction. The kernel detects the planted
pair at $267\times$ the energy of the permitted-but-absent pair, recovers the
component shape at Spearman $0.98$ and localizes where it fires at $0.98$; the
absent pair's field is flat ($3\times10^{-4}$ versus $6.6\times10^{-2}$).
Figure~\ref{fig:kernel} shows the true interaction, the estimated component and
the kernel diagonal share the same corner.

\begin{table}[ht]
\centering
\caption{Simulated kernel view: planted interaction on $(u_1,u_2)$, absent on $(u_3,u_4)$.}
\label{tab:sim-kernel}
\begin{tabular}{lr}
\toprule
Quantity & Value\\
\midrule
detection: $A_{12}$ (planted) / $A_{34}$ (absent) & $267\times$\\
recovery: Spearman(estimated component, true interaction) & $+0.983$\\
localization: Spearman(kernel diagonal, true interaction energy) & $+0.984$\\
absent-pair field: mean $K_{34}(x,x)$ / mean $K_{12}(x,x)$ & $0.005$\\
\bottomrule
\end{tabular}
\end{table}

\begin{figure}[ht]
\centering
\includegraphics[width=\textwidth]{artifacts/kernel_view_test.png}
\caption{True purified interaction, estimated component and kernel diagonal
$K_{12}(x,x)$ over $(u_1,u_2)$. The kernel diagonal is a local interaction-energy field.}
\label{fig:kernel}
\end{figure}

\subsection{Simulated: depth and order}

We plant a pure two-way interaction on $(u_1,u_2)$ and a pure three-way
interaction on $(u_3,u_4,u_5)$, both centered products with no lower-order
content, and fit unconstrained models at depth $1$ to $4$.
Table~\ref{tab:sim-depth} shows each order is recovered when depth reaches it: the
two-way recovery jumps from $-0.71$ at depth $1$ to $+0.97$ at depth $2$, and the
three-way from $+0.42$ at depth $2$ to $+0.93$ at depth $3$. A depth-two model
cannot represent the three-way term at any number of trees; the recovery stays
flat until depth three. Held-out error falls with depth and flattens after depth
three, where no higher order remains (Figure~\ref{fig:depth}).

\begin{table}[ht]
\centering
\caption{Simulated depth and order. Recovery is Spearman of the model's order-$m$
contrast against the planted order-$m$ component.}
\label{tab:sim-depth}
\begin{tabular}{lrrr}
\toprule
Depth & Held-out MSE & Recover 2-way & Recover 3-way\\
\midrule
1 & 0.341 & $-0.712$ & $+0.467$\\
2 & 0.077 & $+0.974$ & $+0.418$\\
3 & 0.045 & $+0.983$ & $+0.928$\\
4 & 0.036 & $+0.986$ & $+0.941$\\
\bottomrule
\end{tabular}
\end{table}

\begin{figure}[ht]
\centering
\includegraphics[width=\textwidth]{artifacts/depth_order_test.png}
\caption{Left: held-out error versus depth. Right: component recovery versus depth;
the two-way effect turns on at depth two, the three-way at depth three.}
\label{fig:depth}
\end{figure}

\subsection{California housing}

The interaction off-diagonal (shared by both matrices) is dominated by
Latitude$\times$Longitude ($A=0.0082$), followed by MedInc$\times$Latitude
($0.0028$) and MedInc$\times$Longitude ($0.0023$): the income effect varies by
location, and location interacts with itself as a spatial surface. The two
diagonals separate importance from bending. MedInc is the third most important
main effect ($A_j=0.026$) yet nearly linear (self-curvature $0.0005$), while
Latitude and Longitude are both important and curved ($A_j=0.067,0.048$;
self-curvature $0.0049,0.0044$) and AveOccup is weak in importance but relatively
curved ($0.006$; $0.0025$). The interaction field for Latitude$\times$Longitude
(Figure~\ref{fig:cal}, right) concentrates along the coast. The depth-order MSE is
$0.0364,0.0257,0.0220,0.0206,0.0190,0.0183$ for depth $1$ to $6$: the largest drop
is depth one to two ($0.011$) and the curve flattens by depth three, so the
effective interaction order is near two.

\begin{figure}[ht]
\centering
\includegraphics[width=\textwidth]{artifacts/curvature_california_housing.png}
\caption{California housing. Left: importance matrix (diagonal is main-effect
energy $A_j$). Center: Hessian matrix (diagonal is self-curvature $A_{jj}$). Both
share the interaction off-diagonal. Right: local interaction field for the top pair.}
\label{fig:cal}
\end{figure}

\subsection{Bikeshare}

Hour$\times$workingday dominates the interaction off-diagonal ($A=0.153$),
eighteen times the next pair (hour$\times$weekday, $0.008$): the commute shape of
ridership depends on whether the day is a working day. On the diagonal, the
self-curvature is exactly zero for the two-state features (year, holiday,
workingday), because a binary main effect is a level shift and cannot bend; the
discrete Hessian reports that without the normalization artifact a variance-scaled
curvature would produce. Hour is both the most important main effect ($A_j=1.41$)
and the most curved (self-curvature $0.14$, the twin commute peaks); month is more
curved ($0.043$) than important ($0.031$), a strongly nonlinear seasonal cycle;
temperature is important but smooth. The depth-order MSE is
$0.395,0.219,0.117,0.098,0.090,0.087$: large drops at both depth two ($0.176$) and
depth three ($0.102$), then a flat tail, so bikeshare carries three-way
calendar structure and its effective order is near three. The two depth-order
curves separate an order-two dataset from an order-three one.

\begin{figure}[ht]
\centering
\includegraphics[width=\textwidth]{artifacts/curvature_bikeshare.png}
\caption{Bikeshare. Left: importance matrix. Center: Hessian matrix, with
self-curvature zero for the binary features. Right: local interaction field for
hour$\times$workingday.}
\label{fig:bike}
\end{figure}

\section{Discussion}

The framework separates two things depth is often asked to explain. Interaction
order is how many features act jointly, set by depth. Local sharpness is how
abruptly one effect changes, set by splits and rounds at fixed order. A per-point
attempt to predict required depth from local sharpness fails once level and other
covariates are controlled, because sharpness is the wrong axis; the axis depth
controls is order, and order is read from the mixed-difference tensor and its
kernel blocks.

The estimates depend on the background measure, which sets what a main effect and
an interaction mean; under correlated predictors the components are not exactly
orthogonal and the effect-energy matrix inherits that. The interaction field
$K_S(x,x)$ is a within-model quantity and should not be read causally.

\section{Conclusion}

Interactions are the off-diagonal curvature of a fitted model, estimated for
piecewise-constant trees as mixed second differences and read directly from the
model-induced kernel, whose diagonal maps where each interaction fires. Tree depth
sets the order of mixed curvature a model can carry, so the depth at which
held-out error flattens measures a dataset's effective interaction order: near two
for California housing, near three for bikeshare.

\begin{thebibliography}{9}
\bibitem{Sudjianto2025}
A. Sudjianto.
\newblock From Classical ANOVA to Gradient Boosted Trees: A Functional ANOVA
Perspective on Inherently Interpretable Models.
\newblock SSRN Working Paper 5315091, 2025.

\bibitem{Hu2023}
L. Hu, V. N. Nair, A. Sudjianto, A. Zhang, and J. Chen.
\newblock Interpretable Machine Learning based on Functional ANOVA Framework:
Algorithms and Comparisons.
\newblock arXiv:2305.15670, 2023.

\bibitem{SudjiantoZhang2026}
A. Sudjianto and A. Zhang.
\newblock Generalized Nadaraya--Watson Operators for Learned Kernel Geometry.
\newblock SSRN Working Paper 6823639, 2026.

\bibitem{Hooker2007}
G. Hooker.
\newblock Generalized Functional ANOVA Diagnostics for High-Dimensional Functions
of Dependent Variables.
\newblock \emph{Journal of Computational and Graphical Statistics}, 16(3):709--732, 2007.

\bibitem{FriedmanPopescu2008}
J. H. Friedman and B. E. Popescu.
\newblock Predictive Learning via Rule Ensembles.
\newblock \emph{The Annals of Applied Statistics}, 2(3):916--954, 2008.

\bibitem{ChenGuestrin2016}
T. Chen and C. Guestrin.
\newblock XGBoost: A Scalable Tree Boosting System.
\newblock In \emph{Proceedings of KDD}, 2016.
\end{thebibliography}

\end{document}
